\subsection{Eighth letter}

\begin{quotationx}
Letter VIII repeats the same topics discussed in Letter VII. I'm sure that Evola's letters would be more interesting, but they are unavailable, presumably because Guenon's family was not forthcoming about releasing his personal papers. We still see Guenon's often curt and condescending tone in addressing Evola.

Evola still does not grasp the notions of the possible and the real; this will need to be explored. We see that Evola was still interested in certain marginal figures (from Guenon's point of view). Curiously, Evola was interested in Eliphas Levi. Equally curious is Guenon's admission about the existence of Hermetic organizations and his own ``Western'' initiation (Letter VII)\footnote{\url{https://gornahoor.net/?p=4559}}; in East and West, Guenon denied any existence of initiation in the West. My guess is that Evola was fishing for such organizations, but Guenon was not forthcoming on the grounds that (1) he does not give personal advice and (2) his own experience is of no interest to anyone else. Yet, it strains credulity to believe there were only 12 Hermetists in Europe.

Once again, I have omitted discussions about publishing, etc.
\end{quotationx}


13 June 1949

Cairo, Egypt Discussions about publishing, etc. omitted As you imagine, it's been quite some time since I had a chance to read your Revolt against the Modern Word. I will therefore make an effort to reread it when I can find some free time, in order to see if there are some points to make as you requested.\footnote{Evola, at that time, was revising Revolt to be republished in a new edition. More discussions about publishing, etc. omitted.} 

According to what you explained to me this time, it seems that you consider the words ``possible'' and ``real'' in the sense of ``non-manifested'' and ``manifested''; if that were so, one could say that it is merely a question of terminology and that, in spite of this expressive difference, we are basically in agreement on the point in question. However, such a use of the words ``possible'' and ``real'', in a sense much different from how we use it, does not seem to be acceptable, because the non-manifested is not only just as real, but even more real than the manifested.

What I said last time regarding my ties with initiatic organizations (even though I don't really like to speak of these things that ultimately can be of interest to no one outside of myself) was in response to what you wrote: ``most often out of that secret society those capable of greater comprehension with respect to initiatic things were found, something that perhaps was verified in your own situation.''

That made me think you gave yourself the idea that, in my case, it could be a question of one of those pretended initiations without any regular ties, which, in my opinion, I could consider only purely imaginary. By the way, I will point out to you that, in Perspectives, I dedicated an entire chapter to explain the reasons why the word ``secret society'' is absolutely unacceptable in cases of the type of those which you referred to.

You think that, in Perspectives, we do not speak of Christian Hermetic organizations; but to the contrary, I expressly mentioned them even in the note to which you referenced and, if I didn't talk about it more, it is because those whose existence I was able to come to know admits such a restricted number of members that they can be considered as inaccessible for all practical purposes. I also see that you have not well understood in what sense I spoke of ``complex problems''. I only wanted to say with what in reality they have many more elements than what can be known through a study made ``from the outside''; it is therefore totally contrary to something that could be defined as you thought.

As for the source itself of the question concerning Masonry, I clearly mean that I do not at all claim to convince you, and that otherwise you would have no interest in it. You say that in that case, it is a question for you only of the truth, but it is also the same even for me. You know moreover that I have never been concerned to entice anyone to join one or another organization, no more than to distance him from them. I only said in a very clear way that that could not be my role. I never had the time nor the interest to be concerned with individual cases and I always refused to give particular advise to anyone, for this thing as for any other. That said, I must however make two or three observations on what you tell me this time, and first of all on what concerns the side degrees, since the true nature of the relationship between those and Masonry seems to elude you. When I speak of Masonry without further clarifications, it is always about Masonry properly called, including only the three degrees of Apprentice, Fellow Craft, and Master Mason, to which can only be added the English degrees of Mark and Royal Arch, totally unknown in ``continental'' Masonry.

Regarding the many other degrees like those you refer to, it is obvious that internally there are some things of a quite different character, and that the connection which you wanted to establish between the different ``systems'' is completely artificial. I am furthermore less inclined to question what I myself wrote formally in a recent article; but, as that is the way with which all these things ended up by agglomerating themselves around Masonry, they do not form an integral part of it to any qualification and consequently it is not what is in question. Another point on which I would like to bring your attention is that when you say that the Lodges that had not adhered to the ``speculative'' schisms were not able to do anything to stop or rectify its consequences, it seems that you do not take into account things that nevertheless cover a certain importance, like the reestablishment of the degree of Master, totally unknown by those of 1717, or the action of the ``Ancients' Great Lodge'', whose independent existence continued up until 1813. To say so frankly, I have the impression that you always think only of what Masonry became at a certain period in Italy and France, and that you have no idea at all what concerns Anglo-Saxon Masonry.

To move on to other issues, I confess that I do not understand at all what realizations you mean concerning Eliphas Levi; in fact, like his filiation (or rather like drawing inspirations in his writings, since he himself died before that), there was nothing other than the occultist French movement of the end of the XIX century and the beginning of the XX, on whose insignificance I think we find ourselves in agreement.

For Kremmerz, I know well that a very unclear story is concealed underneath, but that it gives rise to many doubts, at least because I was never able to find any proof about the real existence of the organization to which he would have belonged. In any case, even if he had personally received an authentic initiation, that would still demonstrate nothing for the organization he founded, insofar as there were other cases of the same type (e.g., that of Inayat Khan, who belonged to a regular tariqa in India, but whose self-styled ``Order of Sufis'' corresponded absolutely to nothing); everything that I can say, is that his rituals are more or less ``Egyptian'' like those of Cagliostro\footnote{\url{https://gornahoor.net/?p=764}}!

There could nevertheless be certain realizations totally within that circle, as you say, but they do not go beyond the psychic domain, something that entails nothing of the truly initiated. I add that, after Kremmerz' death, the different groups into which his organization divided appear absolutely not to know where to turn. I notice that, concerning Eliphas Levi, I forgot to cite the use of his works by Albert Pike; but in that case it is a matter of an influence exercised (otherwise indirectly) on the interpretation of the side degrees of the Scottish Rite, something that does not proceed even in the direction you have seen.

There are certainly cases in which an influence of the counter-initiation is quite visible, and among them, it is necessary to include those in which traditional information is present in a manner of a willful parody. This is above all Meyrink's case, something that, well intended, does not mean that he was perforce conscious of the influence that was exercised over him. Here is why I am amazed that you seem to have a certain esteem in regard to Meyrink, and all the more so, since he had belonged to Bo Yin Ra's movement, for which you clearly had no regard.\footnote{As long ago as 1924, Evola reviewed some of Bo Yin Ra's works. Although Evola did not render a fully positive judgment, he conceded that his doctrine had some interesting points.}

In this regard, it is necessary on the other hand for me to make a rectification: certainly there was in Bo Yin Ra a little bit of charlatanism and mystification, but there was at least still more to him, because he was connected with a very strange organization that had its own headquarters in parts of Turkestan and represented a more or less deviant type of Tantrism. About that, I can certainly be certain (and perhaps I am the only one), at the time when the future Bo Yin Ra was still called Joseph Schneider and studied painting in Paris, some members of the organization in discussion made it known to me one day that he was the only European to be a part of it. Later, I also saw the portrait that Bo Yin Ra had made of his ``Master'' and that it was perfectly recognizable for me; in such occasions, I was able on the other hand to question whether even his most intimate disciples knew absolutely anything at all about that, and I was very aware of letting them in on what I myself knew of it.


\flrightit{Posted on 2012-08-06 by Cologero}

\begin{center}* * *\end{center}

\begin{footnotesize}
\begin{sffamily}
\texttt{Mihai on 2012-08-07 at 04:51 said:}

``Yet, it strains credulity to believe there were only 12 Hermetists in Europe.''

He didn't say that there were only 12 in Europe, but that a particular organization only admitted a maximum of 12 members. Which means (and it is implied) that there were others, but who had a similarly exclusive membership.


\hfill

\end{sffamily}
\end{footnotesize}

